% !PRINTED 401
\weberpage{401}
%<<CONT
der Kettenbruchentwickelung aus einer ungeraden Gliederzahl
besteht. Nun ist eine reducirte Zahl $\omega = \{a, b, c\}$ die erste
Wurzel von
\[
(a, \tfrac{1}{2}b, -c) \quad \text{oder} \quad (2a, b, -2c)
\]
(je nachdem $D$ gerade oder ungerade ist) und die zweite Wurzel von
\[
(-a, \tfrac{1}{2}b, c) \quad \text{oder} \quad (-2a, b, 2c).
\]
Sind diese beiden Formen äquivalent, bestehen also die
Perioden von $\omega$ aus einer ungeraden Gliederzahl, so entspricht
jeder zweiseitigen Formenclasse ein, und jedem Paar entgegengesetzter Classen zwei Systeme äquivalenter Zahlen $\omega$, es ist also
die Classenzahl im \gesperrt{Gauss}'schen Sinne ebenso gross, wie die
Anzahl der Perioden der reducirten Zahlen $\omega$.

Wenn aber die Zahl der Periodenglieder gerade ist, so sind
die Formen $(A, B, C), (-A, B, -C)$ nicht äquivalent und die
\gesperrt{Gauss}'sche Classenzahl ist doppelt so gross, als die Zahl der
Perioden der $\omega$.

\begin{center}
\S.\,129.

\gesperrt{Genäherte Berechnung der reellen Wurzeln einer
numerischen Gleichung durch Kettenbrüche.}
\end{center}

Auf die Theorie der Kettenbrüche gründet sich ein von
\gesperrt{Lagrange} herrührendes Verfahren, um die reellen Wurzeln
einer numerischen Gleichung mit beliebiger Annäherung zu berechnen, das sich jetzt mit wenig Worten auseinander setzen
lässt.

Es sei $f(x) = 0$ eine reelle Gleichung, die wenigstens eine
reelle Wurzel hat. Nach einer der verschiedenen Methoden des
VIII. und IX. Abschnittes können wir dann immer eine ganze
Zahl $a_0$ finden, so dass zwischen $a_0$ und $a_0 + 1$ eine oder mehrere
der Wurzeln von $f(x)$ liegen. Transformiren wir also $f(x)$ durch
die Substitution
\[
x = a_0 + \frac{1}{x_1}
\]
in $f_1(x_1)$, so wird $f_1(x_1)$ eine oder mehrere Wurzeln haben, die
grösser als 1 sind. Wir bestimmen also eine \gesperrt{positive ganze}
Zahl $a_1$, so dass zwischen $a_1$ und $a_1 + 1$ eine oder mehrere der
Wurzeln von $f_1(x_1)$ liegen, und transformiren $f_1(x_1)$ durch die
Substitution
%CONT>>
